%% kapa: exam=2 round=19 year=2008 period=3 subject=compensation_law subject_ko=감정평가 및 보상법규 seq=1 number=1 points=40
%% source: Q-Net 감정평가사 기출문제 (cst00309) artlSeq=5024177 file=extracted/1008/5024177/2008년 제19회 감정평가사 2차시험 기출문제/3교시.hwp
%% conversion: kordoc 4.24.51 -> tools/gen2.py (auto); review=proofread-sonnet; flags=none
\begin{longquestion}{【문제 1】}
서울특별시장은 도시관리계획결정에서 정해진 바에 따라 근린공원을 조성하기 위하여 그 사업에 필요한 토지들을 공익사업을 위한 토지 등의 취득 및 보상에 관한 법률의 규정에 의거하여 협의를 거쳐 취득하고자 하였으나 협의가 성립되지 않아 중앙토지수용위원회에 재결을 신청하였다. 중앙토지수용위원회의 수용재결(수용의 개시일: 2005. 6. 30.)에 따라 서울특별시는 보상금을 지급하고 필요한 토지에 대한 소유권이전등기를 마쳤다. 서울특별시장은 토지를 취득한 후, 6개월간의 공사 끝에 공원을 조성하였다. 공원조성공사가 완료된 후 2년이 지난 뒤에 위 토지를 포함한 그 일대의 토지들이 택지개발예정지구로 지정되었다(고시일: 2008. 6. 30.). 국토해양부장관에 의하여 택지개발사업의 시행자로 지정된 대한주택공사는 택지개발사업실시계획의 승인을 얻어 공원시설을 철거하고, 그 지상에 임대주택을 건설하는 공사를 시행하고 있다. 이에 공원조성사업을 위해 수용된 토지의 소유자 甲은 2008. 8. 30. 서울특별시에 환매의 의사표시를 하였으나, 서울특별시는 甲에게 환매권이 없다고 하여 수용된 토지를 되돌려 주지 않았다. 이러한 경우에 甲이 소유권 회복을 위해 제기할 수 있는 소송수단 및 그 인용가능성에 대하여 검토하시오. (40점) \par
\textbf{[참조 조문]} \par
\textbf{공익사업을 위한 토지 등의 취득 및 보상에 관한 법률} \par
제4조 (공익사업) 이 법에 의하여 토지등을 취득 또는 사용할 수 있는 사업은 다음 각호의 1에 해당하는 사업이어야 한다. \par
1. (생략) \par
2. (생략) \par
3. 국가 또는 지방자치단체가 설치하는 청사ㆍ공장ㆍ연구소ㆍ시험소ㆍ보건 또는 문화시설ㆍ공원ㆍ수목원ㆍ광장ㆍ운동장ㆍ시장ㆍ묘지ㆍ화장장ㆍ도축장 그 밖의 공공용 시설에 관한 사업 \par
4. (생략) \par
5. 국가ㆍ지방자치단체ㆍ정부투자기관ㆍ지방공기업 또는 국가나 지방자치단체가 지정한 자가 임대나 양도의 목적으로 시행하는 주택의 건설 또는 택지의 조성에 관한 사업 \par
6. (생략) \par
7. (생략) \par
8. (생략) \par
\end{longquestion}
